\documentclass{article}
\usepackage{amsmath,amssymb}
\title{Four-direction Wilson transfer dissemination}\author{}\date{}
\begin{document}\maketitle
The argument is valid for all integer side lengths $n\geq3$. It needs no commutation of the modified and original transfer operators, no parity condition, and no change of torus. Coupling conventions are fixed explicitly below.

Use $s(U)=1-N^{-1}\operatorname{Re}\operatorname{tr}U$ and coefficient $\beta_W$. The operator note instead uses $s_{\rm AI}=Ns$ and coefficient $\beta_{\rm AI}=\beta_W/N$. Thus its $B=N\beta_{\rm AI}$ is precisely $\beta_W$ here.

\section{Positive operator and exact original normalization}
Fix one direction as time. On spatial link configurations $X$ use normalized product Haar measure. Let $P$ average the site gauge action $(gX)_e=g_xX_eg_y^{-1}$, and let
\[
 Q_{\boldsymbol\lambda}(X,Y)=\prod_e
 e^{-\lambda_e[N-\operatorname{Re}\operatorname{tr}(X_eY_e^{-1})]}.
\]
Each one-link central convolution is positive: expanding the exponential into powers of the fundamental character and its conjugate gives nonnegative sums of tensor-product characters. Its spectral trace is the kernel value at identity, namely one. Hence the finite product $Q$ is positive trace class. For arbitrary link-dependent $\lambda_e\geq0$,
\[
 Q(gX,gY)=Q(X,Y),
\]
because $X_eY_e^{-1}$ is conjugated by $g_x$. Therefore $Q$ commutes with the unitary gauge action and with $P$. The spatial potential multiplier $M(X)=e^{-\beta_W V(X)/2}$ is gauge invariant. Thus $T=MQPM$ is positive trace class.

For $\lambda_e=\beta_W/N$, its kernel can be written exactly as
\[
 T(X,Y)=M(X)M(Y)\int\prod_xdV_x
 \prod_{e=(x,y)}e^{-\beta_W s(X_eV_yY_e^{-1}V_x^{-1})}.
 \tag{1}
\]
In $\operatorname{Tr}T^n$, the $n$ copies of $X$ are all original spatial links, and the $n$ sets of $V_x$ are all original temporal links. Neighboring half-potentials supply each spatial plaquette once. Periodic trace closure preserves the temporal holonomy integration. Every Haar measure is normalized. Thus $\operatorname{Tr}T^n=Z_{n,\beta_W}$ without a scalar gauge-volume factor. One may algebraically combine the commuting projections, but that simplification is not needed for this identification.

\section{The two inequalities and noncommutativity}
For positive trace-class $A,B$, Schatten Holder gives
\[
 \operatorname{Tr}(AB^{n-1})\leq
 (\operatorname{Tr}A^n)^{1/n}(\operatorname{Tr}B^n)^{(n-1)/n}.
\]
It can also be proved in an eigenbasis of $B$ by scalar Young followed by
$\langle i,Ai\rangle^n\leq\langle i,A^ni\rangle$. No $AB=BA$ is assumed.

For a positive bounded, gauge-invariant slice insertion $D$, take
$A=T^{1/2}DT^{1/2}$ and $B=T$. Cyclicity identifies the left side with $\operatorname{Tr}(DT^n)$. The operators $T^{1/2}DT^{1/2}$ and $D^{1/2}TD^{1/2}$ have the same nonzero eigenvalues (they are $C^*C$ and $CC^*$), so the right-side trace is that of a repeated modified spatial potential. The boundedness of $D$ holds at each finite torus because the Wilson action is bounded.

For a marked temporal slab, the insertion decreases selected $\lambda_e$ from $\beta_W/N$ to $(1-t)\beta_W/N$. The modified $T_1$ remains positive and trace class; the original marked numerator is exactly $\operatorname{Tr}(T_1T^{n-1})$. Holder repeats that same slab. It is not necessary that $T_1$ commute with $T$.

\section{Explicit placement in four directions}
For a plaquette of orientation $(1,2)$, let $A_j$ be marked base points. They progress as follows:
\[
\begin{array}{c|c|c}
 j&\text{free coordinates of the marked bases}&\text{next transfer direction/type}\\\hline
0&\varnothing&3,\ \text{spatial slice}\\
1&3&4,\ \text{spatial slice}\\
2&3,4&1,\ \text{temporal slab}\\
3&1,3,4&2,\ \text{temporal slab}\\
4&1,2,3,4&\text{finished}
\end{array}
\]
In the third step only direction-2 spatial links with the chosen coordinate-2 base are modified. In the fourth step all direction-1 links of the chosen time-2 slab are modified. Reversing the orientation convention for a time-2 plaquette leaves $s(U^{-1})=s(U)$ unchanged.

Writing $F(A)$ for the marked numerator divided by the same original uniform partition function, each step gives $F(A_j)\leq F(A_{j+1})^{1/n}$. Thus
\[
 E_{n,\beta_W}e^{t\beta_Ws(U_P)}
 \leq\left[\frac{Z_{n,(1-t)\beta_W}}{Z_{n,\beta_W}}\right]^{1/n^4},
 \qquad 0\leq t<1.
 \tag{2}
\]
The final anisotropic numerator is no larger than the fully reduced-coupling numerator because every remaining plaquette action is nonnegative. The four applications never change the denominator measure.

Using the actual tree/chain partition exponents $Q_n=3n^4+1$ and $R_n=3n^4-4n^3+1$ gives
\[
 E e^{\beta_Ws(U_P)/2}\leq C_G\beta_W^{2d_G/n},\quad\beta_W\geq1.
 \tag{3}
\]
This is the same original Wilson measure on odd or even tori. Uniformity of the local exponential moment itself requires $\log\beta_W/n$ bounded; the displayed finite-size factor is retained otherwise.

\section{Consequence for the complete-cube result}
Replace (1) of the averaged complete cube calculation by (3). Its exact bath-event proof is unchanged and now gives
\[
 \Pr(\mathcal G^c)\leq864C_G\beta_W^{2d_G/n}
 e^{-\beta_W\varepsilon_0/2}
\]
for every integer $n$ large enough to embed the box, without a parity restriction. Therefore its averaged conditional approximation and raw bad-log mass estimates apply to the original odd-$L$ physical sequence when the records fit. This independent derivation agrees with the operator note's four-direction argument. It does not prove a multilocal dissemination estimate for arbitrary marked sets, which would require a separate placement argument.
\end{document}
